% Source: Gerzensee "Real Analysis" slides 49-57 (epsilon-delta continuity, the
% sin(x)/x and sin(1/x) examples, continuity of linear functionals and of the
% norm, separate versus joint continuity, composition, max of finitely many
% continuous functions, sequential characterization).
% Supplementary: Fukuda (2026), Sec. 7.4.2 (pp. 239-252), Sec. 7.5 (pp. 254-256),
% Sec. 8.1 (compactness and continuity, pp. 261-270), Sec. 8.3 (semicontinuity
% and the Weierstrass theorem, pp. 271-278).
% External references actually cited in this chapter: Rudin (1976).

\chapter{Continuity, Semicontinuity, and the Weierstrass Theorem}
\label{ch:continuity}

\section{Motivation: Why Optima Exist}\label{sec:cont-motivation}

An economic model that begins ``the consumer chooses $x$ to maximise $u(x)$
subject to $\ip{p}{x}\leq m$'' has presupposed that a maximiser exists. Nothing
guarantees this. The supremum of a function over a set need not be attained, and
when it is not, the demand function is undefined, comparative statics are
vacuous, and equilibrium existence arguments have no starting point.

The Weierstrass theorem gives the standard sufficient conditions: a compact
feasible set and enough continuity of the objective. \autoref{ch:topology}
supplied compactness. This chapter supplies continuity, and then weakens it as
far as the existence conclusion permits --- to upper semicontinuity for maxima
and lower semicontinuity for minima. The weakening is not decoration. Indirect
utility functions, value functions of dynamic programs and profit functions
arise as suprema of families of continuous functions, and
\autoref{prop:sup-semicontinuous} shows that such suprema are in general only
semicontinuous.

\section{Continuity}\label{sec:cont-basic}

Throughout, $A\subseteq\R^n$ and $f\colon A\to\R^m$ unless stated otherwise. Two
metrics are in play --- the Euclidean metric on the domain and on the codomain
--- and both are written $d$.

\begin{definition}[Relative topology]\label{def:relative-topology}
	Let $A\subseteq X$ with $(X,d)$ a metric space. A set $U\subseteq A$ is
	\dfnas{open relative to $A$}{relatively open set} if $U=G\cap A$ for some open
	$G\subseteq X$; equivalently, if for every $x\in U$ there is $\varepsilon>0$
	with $\ball{\varepsilon}{x}\cap A\subseteq U$. Relative closedness is defined
	analogously. This is exactly the topology induced by restricting $d$ to
	$A\times A$, so ``relatively open in $A$'' and ``open in the metric space
	$(A,\restr{d}{A\times A})$'' are the same statement.
\end{definition}

The distinction matters because $[0,1)$ is not open in $\R$ but is open relative
to $[0,1]$, and hypotheses about feasible sets are almost always relative
hypotheses.

\begin{definition}[Continuity]\label{def:continuity}
	Let $A\subseteq\R^n$ and $f\colon A\to\R^m$. The mapping $f$ is
	\dfnas{continuous at}{continuity} $x_0\in A$ if for every $\varepsilon>0$ there
	is $\delta>0$ such that
	\[
		x\in A\ \text{and}\ d(x_0,x)<\delta
		\quad\Longrightarrow\quad
		d\bigl(f(x_0),f(x)\bigr)<\varepsilon .
	\]
	It is \dfnas{continuous on $A$}{continuous mapping} if it is continuous at every
	point of $A$.
\end{definition}

Two features of the definition are easy to pass over. First, $\delta$ may depend
on $x_0$ as well as on $\varepsilon$; requiring that it not is uniform
continuity (\autoref{def:uniform-continuity}). Second, the implication is
restricted to $x\in A$, so continuity is a statement relative to the domain. A
function defined only on $\Z$ is continuous at every point of its domain, for
trivial reasons.

\begin{eg}[Removable and essential discontinuities]\label{eg:sinx-over-x}
	Define $f\colon\R\to\R$ by $f(x)=\sin(x)/x$ for $x\neq0$ and $f(0)=1$. Then $f$
	is continuous at $0$: since $\cos x\leq\sin(x)/x\leq1$ for $0<\abs{x}<\pi/2$ and
	$\cos$ is continuous with $\cos0=1$, the value $1$ is the only value at $0$ that
	makes $f$ continuous, and it has been assigned. The discontinuity of the
	function $x\mapsto\sin(x)/x$ on $\R\setminus\{0\}$ is \emph{removable}.

	Define $g\colon\R\to\R$ by $g(x)=\sin(1/x)$ for $x\neq0$ and $g(0)=0$. No
	assignment at $0$ makes $g$ continuous. Taking $x_k=\bigl(\pi/2+2k\pi\bigr)^{-1}$
	and $y_k=\bigl(-\pi/2+2k\pi\bigr)^{-1}$ gives $x_k\to0$, $y_k\to0$ but
	$g(x_k)=1$, $g(y_k)=-1$, so no value of $g(0)$ is compatible with
	\autoref{prop:sequential-continuity}. The discontinuity is \emph{essential}
	\autocite[slides 50--51]{fukuda2026analysis}.
\end{eg}

\begin{proposition}[Continuity of linear functionals and of the norm]
	\label{prop:linear-continuous}
	Let $a\in\R^n$.
	\begin{enumerate}[(i)]
		\item The functional $f(x)=\ip{a}{x}$ is continuous on $\R^n$.
		\item Every linear map $T\colon\R^n\to\R^m$ is continuous.
		\item The norm $\norm{\cdot}\colon\R^n\to\R$ is continuous.
	\end{enumerate}
\end{proposition}

\begin{proof}
	(i) If $a=0$ then $f\equiv0$ is constant, hence continuous. Let $a\neq0$, fix
	$x_0\in\R^n$ and $\varepsilon>0$, and set $\delta\eqdef\varepsilon/\norm{a}>0$.
	If $\norm{x-x_0}<\delta$ then by Cauchy--Schwarz
	(\autoref{lem:cauchy-schwarz})
	\[
		\abs{f(x)-f(x_0)}=\abs{\ip{a}{x-x_0}}\leq\norm{a}\norm{x-x_0}
		<\norm{a}\delta=\varepsilon .
	\]

	(ii) Write $T$ in the canonical bases as in \autoref{thm:matrix-rep}, so that the
	$i$th component of $Tx$ is $\ip{a_i}{x}$ with $a_i$ the $i$th row of the
	representing matrix. Given $\varepsilon>0$, apply (i) to each component with
	tolerance $\varepsilon/\sqrt m$ and take the smallest of the resulting $\delta_i$.

	(iii) By the reverse triangle inequality (\autoref{prop:reverse-triangle}),
	$\abs{\norm{x}-\norm{x_0}}\leq\norm{x-x_0}$, so $\delta=\varepsilon$ works.
\end{proof}

Part (i) is Proposition~7.17 of \textcite{fukuda2026notes}, proved on
\autocite[slide 52]{fukuda2026analysis}; part (iii) is Proposition~7.18 there,
whose hint is Exercise~7.4, recorded here as
\autoref{prop:reverse-triangle}. The $\delta$ produced in (i) does not depend on
$x_0$, so linear functionals are in fact uniformly continuous; this is a
consequence of linearity, not of the specific bound.

\begin{proposition}[Algebra of continuous functions]\label{prop:continuity-algebra}
	Let $A\subseteq\R^n$ and let $f,g\colon A\to\R$ be continuous at $x_0$.
	\begin{enumerate}[(i)]
		\item $f+g$, $fg$ and $\alpha f$ for $\alpha\in\R$ are continuous at $x_0$;
		      $f/g$ is continuous at $x_0$ provided $g(x_0)\neq0$.
		\item If $h\colon B\to\R^m$ is continuous at $f(x_0)$ and $f(A)\subseteq B$,
		      then $h\circ f$ is continuous at $x_0$.
		\item If $f_1,\dots,f_m\colon A\to\R$ are continuous at $x_0$, then so are
		      $\max_{1\leq i\leq m}f_i$ and $\min_{1\leq i\leq m}f_i$.
	\end{enumerate}
\end{proposition}

\begin{proof}
	(i) and (ii) are the standard $\varepsilon$--$\delta$ arguments; see
	\textcite[Thms.~4.4 and~4.7]{rudin1976principles}.

	(iii) For $m=2$ use the identity
	$\max\{s,t\}=\tfrac12\bigl(s+t+\abs{s-t}\bigr)$ and combine (i), (ii) and the
	continuity of $\abs{\cdot}$; the general case follows by induction on $m$, since
	$\max_{i\leq m}f_i=\max\{\max_{i\leq m-1}f_i,f_m\}$. The corresponding identity
	for the minimum is $\min\{s,t\}=\tfrac12(s+t-\abs{s-t})$.
\end{proof}

Part (iii) is Exercise~7.34 of \autocite[slide 56]{fukuda2026analysis}. The
restriction to \emph{finitely} many functions is essential and is exactly what
fails in \autoref{eg:sup-not-usc}; the induction consumes one step per function
and cannot be run to infinity.

\begin{proposition}[Sequential characterisation]\label{prop:sequential-continuity}
	Let $A\subseteq\R^n$ and $f\colon A\to\R^m$. Then $f$ is continuous at $x_0\in A$
	if and only if for every sequence $(x_k)\subseteq A$ with $x_k\to x_0$ one has
	$f(x_k)\to f(x_0)$; that is,
	$f\bigl(\lim_{k}x_k\bigr)=\lim_{k}f(x_k)$.
\end{proposition}

\begin{proof}
	($\Rightarrow$) Given $\varepsilon>0$, take $\delta$ from
	\autoref{def:continuity} and then $k_0$ with $d(x_0,x_k)<\delta$ for $k\geq k_0$;
	then $d(f(x_0),f(x_k))<\varepsilon$ for $k\geq k_0$.

	($\Leftarrow$) Suppose $f$ is not continuous at $x_0$. Then there is
	$\varepsilon_0>0$ such that for every $\delta>0$ some $x\in A$ has
	$d(x_0,x)<\delta$ and $d(f(x_0),f(x))\geq\varepsilon_0$. Taking
	$\delta=1/k$ produces $x_k\in A$ with $x_k\to x_0$ and
	$d(f(x_0),f(x_k))\geq\varepsilon_0$ for every $k$, so $f(x_k)\not\to f(x_0)$.
\end{proof}

This is Proposition~7.22 of \textcite{fukuda2026notes}
\autocite[slide 57]{fukuda2026analysis}. The contrapositive is the practical
tool: to refute continuity, exhibit two sequences approaching the same point
along which $f$ has different limits, as in
\autoref{eg:separate-vs-joint}.

\begin{theorem}[Topological characterisation]\label{thm:continuity-preimage}
	Let $A\subseteq\R^n$ and $f\colon A\to\R^m$. The following are equivalent.
	\begin{enumerate}[(i)]
		\item $f$ is continuous on $A$.
		\item $f^{-1}(V)$ is open relative to $A$ for every open $V\subseteq\R^m$.
		\item $f^{-1}(F)$ is closed relative to $A$ for every closed $F\subseteq\R^m$.
	\end{enumerate}
\end{theorem}

\begin{proof}
	(i)$\Rightarrow$(ii). Let $V$ be open and $x_0\in f^{-1}(V)$. Since $V$ is open
	there is $\varepsilon>0$ with $\ball{\varepsilon}{f(x_0)}\subseteq V$;
	continuity supplies $\delta>0$ with
	$f\bigl(\ball{\delta}{x_0}\cap A\bigr)\subseteq\ball{\varepsilon}{f(x_0)}
		\subseteq V$, that is $\ball{\delta}{x_0}\cap A\subseteq f^{-1}(V)$.

	(ii)$\Rightarrow$(i). Fix $x_0\in A$ and $\varepsilon>0$. The set
	$V=\ball{\varepsilon}{f(x_0)}$ is open, so $f^{-1}(V)$ is relatively open and
	contains $x_0$; hence some $\delta>0$ has
	$\ball{\delta}{x_0}\cap A\subseteq f^{-1}(V)$, which is the required
	implication.

	(ii)$\Leftrightarrow$(iii). Preimages commute with complements:
	$f^{-1}(\comp{F})=\comp{\bigl(f^{-1}(F)\bigr)}$ relative to $A$
	(\autoref{prop:image-preimage}). Apply
	\autoref{prop:interior-closure-duality}.
\end{proof}

\begin{remark}[Continuity, measurability, and the shape of the definition]
	\label{rem:continuity-measurability}
	\autoref{thm:continuity-preimage} states continuity without reference to any
	metric: only the two collections of open sets are used. This is the definition of
	continuity in a general topological space, and it is what makes the parallel of
	\autoref{rem:topology-sigma-algebra} exact. A map is continuous when preimages
	of open sets are open; a map is measurable when preimages of measurable sets are
	measurable (\autoref{def:random-variable}). The same asymmetry appears in both:
	the condition constrains preimages, never images. Images of open sets under
	continuous maps need not be open --- $x\mapsto x^2$ sends the open set
	$(-1,1)$ to $[0,1)$ --- and this is why \autoref{prop:continuous-compact}, which
	does concern images, requires compactness rather than openness.
\end{remark}

\begin{eg}[Separate continuity does not imply joint continuity]
	\label{eg:separate-vs-joint}
	Let $f\colon\R^p\times\R^q\to\R^m$ be continuous. For fixed $y_0$ the section
	$x\mapsto f(x,y_0)$ is continuous, since $x\mapsto(x,y_0)$ is continuous and
	\autoref{prop:continuity-algebra}(ii) applies; symmetrically in the other
	argument. The converse fails. Define $f\colon\R^2\to\R$ by
	\[
		f(x_1,x_2)=
		\begin{cases}
			\dfrac{x_1x_2}{x_1^2+x_2^2}, & (x_1,x_2)\neq(0,0), \\[1.2ex]
			0,                           & (x_1,x_2)=(0,0).
		\end{cases}
	\]
	Both sections through the origin are identically zero, hence continuous. But
	along the diagonal $f(\rho,\rho)=\tfrac12$ and along the anti-diagonal
	$f(\rho,-\rho)=-\tfrac12$ for every $\rho\neq0$, so by
	\autoref{prop:sequential-continuity} $f$ is not continuous at the origin. In
	fact $f$ is constant on every ray through the origin, and takes every value in
	$[-\tfrac12,\tfrac12]$ in every ball around it
	\autocite[slides 54--55]{fukuda2026analysis}.

	The example is a standing warning in comparative statics. Verifying that a
	function of prices and income responds continuously to each argument separately
	establishes nothing about its joint behaviour, and it is joint continuity that
	the maximum theorem of \autoref{ch:maximum} requires.
\end{eg}

\subsection{Continuity, compactness, and connectedness}

\begin{proposition}[Continuous images of compact sets]\label{prop:continuous-compact}
	Let $K\subseteq\R^n$ be compact and $f\colon K\to\R^m$ continuous. Then $f(K)$
	is compact.
\end{proposition}

\begin{proof}
	Let $(y_k)\subseteq f(K)$ and choose $x_k\in K$ with $f(x_k)=y_k$. By
	compactness of $K$ some subsequence $x_{k_j}\to x^\ast\in K$, and by
	\autoref{prop:sequential-continuity} $y_{k_j}=f(x_{k_j})\to f(x^\ast)\in f(K)$.
	So $f(K)$ is sequentially compact, hence compact by
	\autoref{thm:compact-seq-compact}.
\end{proof}

\begin{proposition}[Continuous images of connected sets]\label{prop:continuous-connected}
	Let $C\subseteq\R^n$ be connected and $f\colon C\to\R^m$ continuous. Then $f(C)$
	is connected.
\end{proposition}

\begin{proof}
	Suppose $f(C)=U_1\cup U_2$ with $U_1,U_2$ non-empty, disjoint and open relative
	to $f(C)$. By \autoref{thm:continuity-preimage} the sets $f^{-1}(U_1)$ and
	$f^{-1}(U_2)$ are open relative to $C$, non-empty, disjoint, and their union is
	$C$ --- contradicting connectedness of $C$.
\end{proof}

\begin{corollary}[Intermediate value theorem]\label{cor:ivt}
	Let $f\colon[a,b]\to\R$ be continuous and let $y$ lie between $f(a)$ and $f(b)$.
	Then $f(c)=y$ for some $c\in[a,b]$.
\end{corollary}

\begin{proof}
	$[a,b]$ is connected by \autoref{prop:connected-interval}, so $f([a,b])$ is
	connected by \autoref{prop:continuous-connected}, hence an interval by
	\autoref{prop:connected-interval} again. An interval containing $f(a)$ and
	$f(b)$ contains everything between them.
\end{proof}

\subsection{Uniform and Lipschitz continuity}

\begin{definition}[Uniform, Lipschitz and H\"older continuity]\label{def:uniform-continuity}
	Let $A\subseteq\R^n$ and $f\colon A\to\R^m$.
	\begin{enumerate}[(i)]
		\item $f$ is \dfnas{uniformly continuous}{uniform continuity} on $A$ if for
		      every $\varepsilon>0$ there is $\delta>0$, \emph{depending only on
			      $\varepsilon$}, such that $x,y\in A$ and $d(x,y)<\delta$ imply
		      $d(f(x),f(y))<\varepsilon$.
		\item $f$ is \dfnas{Lipschitz continuous}{Lipschitz continuity} with constant
		      $L\geq0$ if $d(f(x),f(y))\leq L\,d(x,y)$ for all $x,y\in A$. If
		      $L<1$ the map is a contraction (\autoref{def:contraction}).
		\item $f$ is \dfnas{H\"older continuous}{Holder continuity@H\"older continuity}
		      of order $\alpha\in(0,1]$ if $d(f(x),f(y))\leq L\,d(x,y)^\alpha$ for all
		      $x,y\in A$.
	\end{enumerate}
	Lipschitz continuity is H\"older continuity of order $1$; H\"older continuity of
	any order implies uniform continuity, with
	$\delta=(\varepsilon/L)^{1/\alpha}$; and uniform continuity implies continuity.
	None of the three implications reverses.
\end{definition}

\begin{eg}[The implications are strict]\label{eg:continuity-hierarchy}
	On $(0,1]$, $x\mapsto1/x$ is continuous but not uniformly continuous: taking
	$x_k=1/k$ and $y_k=1/(k+1)$ gives $\abs{x_k-y_k}\to0$ while
	$\abs{f(x_k)-f(y_k)}=1$ for every $k$. On $[0,1]$, $x\mapsto\sqrt x$ is
	uniformly continuous --- it is H\"older of order $\tfrac12$ --- but not
	Lipschitz, since $\bigl(\sqrt x-\sqrt0\bigr)/(x-0)=x^{-1/2}\to\infty$ as
	$x\downarrow0$.
\end{eg}

\begin{theorem}[Heine--Cantor]\label{thm:heine-cantor}
	A continuous function on a compact set is uniformly continuous.
\end{theorem}

\begin{proof}
	Let $K$ be compact, $f$ continuous on $K$, and suppose $f$ is not uniformly
	continuous. Then some $\varepsilon_0>0$ admits, for each $k\in\N$, points
	$x_k,y_k\in K$ with $d(x_k,y_k)<1/k$ and
	$d(f(x_k),f(y_k))\geq\varepsilon_0$. By compactness pass to a subsequence with
	$x_{k_j}\to x^\ast\in K$; then $y_{k_j}\to x^\ast$ as well, because
	$d(y_{k_j},x^\ast)\leq d(y_{k_j},x_{k_j})+d(x_{k_j},x^\ast)\to0$. Continuity and
	\autoref{prop:sequential-continuity} give
	$f(x_{k_j})\to f(x^\ast)$ and $f(y_{k_j})\to f(x^\ast)$, so
	$d(f(x_{k_j}),f(y_{k_j}))\to0$, contradicting the lower bound $\varepsilon_0$.
\end{proof}

Uniform continuity is what allows a single $\delta$ to be quoted for an entire
domain, and it is the hypothesis under which Riemann sums converge in
\autoref{ch:integration} and under which numerical approximation error can be
bounded uniformly over a state space in \autoref{ch:dynamic}.

\section{Semicontinuity}\label{sec:cont-semi}

Continuity is more than existence of optima requires. The following weakening
isolates exactly what is needed, and it is the form in which the hypothesis is
verifiable for the value functions that arise in economics.

\begin{definition}[Semicontinuity]\label{def:semicontinuity}
	Let $A\subseteq\R^n$ and $f\colon A\to\R$.
	\begin{enumerate}[(i)]
		\item $f$ is \dfnas{lower semicontinuous}{lower semicontinuity} (lsc) if
		      $f^{-1}\bigl((\lambda,\infty)\bigr)=\{x\in A : f(x)>\lambda\}$ is open
		      relative to $A$ for every $\lambda\in\R$; equivalently, if every lower
		      section $\{x\in A: f(x)\leq\lambda\}$ is closed relative to $A$.
		\item $f$ is \dfnas{upper semicontinuous}{upper semicontinuity} (usc) if $-f$
		      is lower semicontinuous; equivalently, if every upper section
		      $\{x\in A: f(x)\geq\lambda\}$ is closed relative to $A$.
	\end{enumerate}
	\nidx{lsc}{lsc, lower semicontinuous}\nidx{usc}{usc, upper semicontinuous}
\end{definition}

\begin{proposition}[Pointwise form]\label{prop:semicontinuity-pointwise}
	Let $A\subseteq\R^n$ and $f\colon A\to\R$.
	\begin{enumerate}[(i)]
		\item $f$ is lsc if and only if for every $x_0\in A$ and every
		      $\lambda<f(x_0)$ there is $\delta>0$ such that
		      $x\in\ball{\delta}{x_0}\cap A$ implies $\lambda<f(x)$.
		\item $f$ is usc if and only if for every $x_0\in A$ and every
		      $\lambda>f(x_0)$ there is $\delta>0$ such that
		      $x\in\ball{\delta}{x_0}\cap A$ implies $f(x)<\lambda$.
	\end{enumerate}
	Equivalently, in sequential form: $f$ is lsc at $x_0$ if
	$x_k\to x_0$ implies $f(x_0)\leq\liminf_kf(x_k)$, and usc at $x_0$ if
	$x_k\to x_0$ implies $\limsup_kf(x_k)\leq f(x_0)$.
\end{proposition}

\begin{proof}
	(i) If $f$ is lsc and $\lambda<f(x_0)$, then $x_0$ belongs to the relatively open
	set $\{f>\lambda\}$, which therefore contains
	$\ball{\delta}{x_0}\cap A$ for some $\delta>0$. Conversely, given $\lambda$ and
	$x_0\in\{f>\lambda\}$, the hypothesis produces such a $\delta$, so
	$\{f>\lambda\}$ is relatively open. Part (ii) follows by replacing $f$ with
	$-f$.

	For the sequential form, suppose $f$ is lsc at $x_0$ and $x_k\to x_0$. For any
	$\lambda<f(x_0)$ the ball $\ball{\delta}{x_0}$ eventually contains every $x_k$,
	so $f(x_k)>\lambda$ eventually and $\liminf_kf(x_k)\geq\lambda$; letting
	$\lambda\uparrow f(x_0)$ gives the claim. The converse is the same argument run
	backwards.
\end{proof}

This is Proposition~8.10 of \textcite[p.~271]{fukuda2026notes}. The verbal
statement is worth fixing: at a point of lower semicontinuity the function
cannot \emph{drop} suddenly as one moves away from $x_0$, and at a point of upper
semicontinuity it cannot \emph{jump up}. A jump downward is compatible with
upper semicontinuity, which is why the maximum of a usc function can sit exactly
at the top of a jump.

\begin{corollary}\label{cor:semicontinuity-continuity}
	$f$ is continuous if and only if it is both lower and upper semicontinuous.
\end{corollary}

\begin{proof}
	If $f$ is continuous, then $f^{-1}$ of the open sets $(\lambda,\infty)$ and
	$(-\infty,\lambda)$ are relatively open by \autoref{thm:continuity-preimage},
	giving both properties. Conversely, if both hold then
	$f^{-1}\bigl((\lambda,\mu)\bigr)=f^{-1}((\lambda,\infty))\cap
		f^{-1}((-\infty,\mu))$ is relatively open, and every open subset of $\R$ is a
	countable union of such intervals, so \autoref{thm:topology-axioms}(i) and
	\autoref{thm:continuity-preimage} apply.
\end{proof}

\begin{proposition}[Epigraph and hypograph]\label{prop:epigraph}
	Let $A\subseteq\R^n$ and $f\colon A\to\R$. Then $f$ is lsc if and only if its
	\dfn{epigraph}
	\[
		\epi f\eqdef\{(x,z)\in A\times\R : f(x)\leq z\}
	\]
	is closed in $A\times\R$, and $f$ is usc if and only if its
	\dfn{hypograph} $\hyp f\eqdef\{(x,z)\in A\times\R : f(x)\geq z\}$ is closed.
	\nidx{Epi f}{$\epi f$, epigraph}
\end{proposition}

\begin{proof}
	Suppose $f$ is lsc and $(x_k,z_k)\to(x,z)$ with $f(x_k)\leq z_k$. By
	\autoref{prop:semicontinuity-pointwise} gives
	$f(x)\leq\liminf_kf(x_k)\leq\liminf_kz_k=z$, so that $(x,z)\in\epi f$; now
	apply \autoref{thm:closed-sequential}. Conversely, if $\epi f$ is closed and
	$x_k\to x$, put $\ell\eqdef\liminf_kf(x_k)$ and pass to a subsequence along
	which $f(x_{k_j})\to\ell$. Then $(x_{k_j},f(x_{k_j}))\in\epi f$ converges to
	$(x,\ell)$, so $(x,\ell)\in\epi f$, that is $f(x)\leq\ell$. The hypograph
	statement follows by applying this to $-f$.
\end{proof}

This is Proposition~8.12 of \textcite[p.~273]{fukuda2026notes}. The epigraph
characterisation is the bridge to \autoref{ch:convex}: convexity of $f$ is
convexity of $\epi f$, and lower semicontinuity is closedness of $\epi f$, so
``closed convex function'' means precisely that the epigraph is a closed convex
set.

\begin{proposition}[Suprema and infima]\label{prop:sup-semicontinuous}
	Let $A\subseteq\R^n$ and let $(f_i)_{i\in I}$ be a family of functions
	$A\to\R$ indexed by an arbitrary set $I$.
	\begin{enumerate}[(i)]
		\item If every $f_i$ is lsc and $f\eqdef\sup_{i\in I}f_i$ is finite on $A$,
		      then $f$ is lsc.
		\item If every $f_i$ is usc and $f\eqdef\inf_{i\in I}f_i$ is finite on $A$,
		      then $f$ is usc.
	\end{enumerate}
\end{proposition}

\begin{proof}
	(i) For any $\lambda$,
	$\{x : \sup_if_i(x)>\lambda\}=\bigcup_{i\in I}\{x : f_i(x)>\lambda\}$: the
	supremum exceeds $\lambda$ exactly when some member does. Each set on the right
	is relatively open, and \autoref{thm:topology-axioms}(i) permits arbitrary
	unions. Part (ii) follows by negation.
\end{proof}

\begin{eg}[A supremum of continuous functions that is not usc]\label{eg:sup-not-usc}
	For each $a\in(0,1]$ define the continuous function $f_a\colon[0,1]\to\R$ by
	$f_a(x)=x/a$ for $0\leq x\leq a$ and $f_a(x)=1$ for $a\leq x\leq1$. Then
	\[
		f(x)\eqdef\sup_{a\in(0,1]}f_a(x)=
		\begin{cases}
			0, & x=0,          \\
			1, & 0<x\leq1 .
		\end{cases}
	\]
	Indeed $f_a(0)=0$ for every $a$, while for $x>0$ the choice $a=x$ gives
	$f_x(x)=1$. The limit function is lsc --- consistent with
	\autoref{prop:sup-semicontinuous}(i) --- and it is not usc, since
	$\{f\geq1\}=(0,1]$ is not closed in $[0,1]$. This is Example~8.4 of
	\textcite[p.~272]{fukuda2026notes}.

	The pattern recurs throughout economics. Indirect utility, the profit function
	and the value function of a dynamic program are all defined as suprema over a
	feasible set indexed by a parameter, and the general theory delivers only one of
	the two semicontinuity properties for free. Establishing the other requires
	additional structure --- compactness of the feasible correspondence and
	continuity of the objective, which is exactly the content of Berge's theorem in
	\autoref{thm:berge-maximum}.
\end{eg}

\section{The Weierstrass Theorem}\label{sec:cont-weierstrass}

\begin{theorem}[Weierstrass]\label{thm:weierstrass}
	Let $A\subseteq\R^n$ be compact and non-empty.
	\begin{enumerate}[(i)]
		\item A lower semicontinuous $f\colon A\to\R$ attains its minimum: there is
		      $x_\ast\in A$ with $f(x_\ast)=\inf_{x\in A}f(x)$, and this infimum is
		      finite.
		\item An upper semicontinuous $f\colon A\to\R$ attains its maximum: there is
		      $x^\ast\in A$ with $f(x^\ast)=\sup_{x\in A}f(x)$, and this supremum is
		      finite.
	\end{enumerate}
\end{theorem}

\begin{proof}
	It suffices to prove (i); (ii) follows by applying (i) to $-f$.

	Write $\alpha\eqdef\inf_{x\in A}f(x)\in[-\infty,\infty)$ and consider the lower
	sections $L_\lambda\eqdef\{x\in A : f(x)\leq\lambda\}$, each closed relative to
	$A$ by \autoref{def:semicontinuity}, hence compact by
	\autoref{prop:compact-stability}(i).

	First, $\alpha>-\infty$. If not, $L_\lambda\neq\emptyset$ for every
	$\lambda\in\R$, and the family $(L_\lambda)_{\lambda\in\R}$ is nested
	decreasing as $\lambda$ decreases, so every finite subfamily has non-empty
	intersection. A nested family of non-empty compact sets has non-empty
	intersection (\autoref{lem:finite-intersection}), so some
	$x\in\bigcap_{\lambda\in\R}L_\lambda$, giving $f(x)\leq\lambda$ for every real
	$\lambda$, which is impossible for a real-valued $f$. Hence $\alpha$ is finite.

	Second, the infimum is attained. The sets $L_\lambda$ for $\lambda>\alpha$ are
	non-empty by definition of the infimum, compact, and nested. By
	\autoref{lem:finite-intersection} again,
	$\bigcap_{\lambda>\alpha}L_\lambda\neq\emptyset$; any $x_\ast$ in this
	intersection satisfies $f(x_\ast)\leq\lambda$ for every $\lambda>\alpha$, hence
	$f(x_\ast)\leq\alpha$, and $f(x_\ast)\geq\alpha$ because $\alpha$ is a lower
	bound. So $f(x_\ast)=\alpha=\min_{x\in A}f(x)$.
\end{proof}

\begin{lemma}[Finite intersection property]\label{lem:finite-intersection}
	Let $(K_i)_{i\in I}$ be a family of non-empty compact subsets of a metric space
	such that every finite subfamily has non-empty intersection. Then
	$\bigcap_{i\in I}K_i\neq\emptyset$.
\end{lemma}

\begin{proof}
	Fix $i_0\in I$ and suppose $\bigcap_{i}K_i=\emptyset$. Then the relatively open
	sets $\comp{K_i}\cap K_{i_0}$, $i\in I$, cover $K_{i_0}$: a point of $K_{i_0}$
	lying in no $\comp{K_i}$ would lie in every $K_i$. By compactness finitely many
	suffice, say for $i_1,\dots,i_m$, which says
	$K_{i_0}\cap K_{i_1}\cap\dots\cap K_{i_m}=\emptyset$ --- contradicting the
	hypothesis on finite subfamilies.
\end{proof}

\begin{corollary}[Weierstrass, continuous case]\label{cor:weierstrass-continuous}
	A continuous real-valued function on a non-empty compact set attains both its
	minimum and its maximum.
\end{corollary}

\begin{proof}
	By \autoref{cor:semicontinuity-continuity} the function is both lsc and usc;
	apply both parts of \autoref{thm:weierstrass}.

	A second proof, independent of semicontinuity, is worth recording because it
	generalises differently. By \autoref{prop:continuous-compact} the image $f(A)$
	is a compact subset of $\R$, hence closed and bounded by
	\autoref{thm:heine-borel}. Boundedness gives finite
	$M\eqdef\sup f(A)$ and $m\eqdef\inf f(A)$, and by
	\autoref{prop:sup-char}(ii) both are points of closure of $f(A)$; closedness
	then gives $M,m\in f(A)$, that is, both are attained.
\end{proof}

This is Corollary~8.6 of \textcite[p.~275]{fukuda2026notes}, with both proofs as
given there.

\begin{eg}[Each hypothesis is needed]\label{eg:weierstrass-failures}
	\begin{enumerate}[(i)]
		\item \emph{Semicontinuity of the wrong type.} Let $f\colon[0,1]\to\R$ be
		      $f(0)=\tfrac12$ and $f(x)=1-x$ for $x\in(0,1]$. This function is lsc, so
		      by \autoref{thm:weierstrass}(i) it attains its minimum, at $x=1$. It has
		      no maximum: $\sup f=1$, approached as $x\downarrow0$, but $f(0)=\tfrac12$.
		      Symmetrically, $g(0)=\tfrac12$ and $g(x)=x$ on $(0,1]$ is usc, attains its
		      maximum at $x=1$, and has no minimum \autocite[Ex.~8.5,
			      p.~274]{fukuda2026notes}.
		\item \emph{Non-compact domain, unbounded.} $f(x)=x$ on $[0,\infty)$ is
		      continuous and attains no maximum.
		\item \emph{Non-compact domain, not closed.} $f(x)=x$ on $[0,1)$ is continuous
		      and bounded, and attains no maximum: the supremum $1$ is not in the image
		      because the point at which it would be attained is not in the domain.
		\item \emph{Infinite dimensions.} On the closed unit ball of the space of
		      \autoref{eg:no-bw-infinite}, the continuous functional
		      $x\mapsto\sum_k2^{-k}x_k$ has supremum not attained on the set of
		      finitely supported points, illustrating that the closed bounded set is not
		      compact and that \autoref{thm:weierstrass} does not apply.
	\end{enumerate}
\end{eg}

\begin{proposition}[The solution set is compact]\label{prop:argmax-compact}
	Let $A\subseteq\R^n$ be compact and non-empty.
	\begin{enumerate}[(i)]
		\item If $f\colon A\to\R$ is usc, then
		      $\argmax_{x\in A}f(x)$ is a non-empty compact subset of $A$.
		\item If $f\colon A\to\R$ is lsc, then $\argmin_{x\in A}f(x)$ is a non-empty
		      compact subset of $A$.
	\end{enumerate}
\end{proposition}

\begin{proof}
	(i) Non-emptiness is \autoref{thm:weierstrass}(ii). Writing
	$M\eqdef\max_{x\in A}f(x)$, the solution set is
	$\{x\in A : f(x)\geq M\}$, which is closed relative to $A$ by upper
	semicontinuity, hence compact by \autoref{prop:compact-stability}(i). Part (ii)
	is the same argument for $-f$.
\end{proof}

This is Proposition~8.13 of \textcite[p.~275]{fukuda2026notes}. It is the first
appearance of the object that \autoref{ch:maximum} studies systematically: the
solution of an optimisation problem is in general a \emph{set}, and the
statement that it varies continuously with the parameters is a statement about
correspondences, not functions. Uniqueness --- which reduces the set to a
singleton and the correspondence to a function --- requires strict quasiconcavity
of $f$ and convexity of $A$, and is supplied in \autoref{prop:strict-quasi-unique}.

\section{Economic Application: Existence of Consumer Demand}\label{sec:cont-demand}

The four questions raised in \autoref{eg:consumer-setup} can now be answered in
part.

\paragraph{Environment.} There are $n$ goods. A consumption bundle is a vector
$x\in\R^n_+$. Prices are $p\in\R^n_{++}$ and wealth is $m>0$. The consumer's
feasible set is the budget set
$\budget(p,m)=\{x\in\R^n_+ : \ip{p}{x}\leq m\}$, and preferences are
represented by a utility function $u\colon\R^n_+\to\R$.

\paragraph{Problem.} The consumer solves
\begin{equation}\label{eq:ump}
	\max_{x}\ u(x)
	\qquad\st\quad \ip{p}{x}\leq m,\quad x\geq0 .
\end{equation}

\begin{proposition}[Existence of a solution to the consumer's problem]
	\label{prop:demand-exists}
	Let $p\in\R^n_{++}$, $m>0$, and let $u\colon\R^n_+\to\R$ be upper
	semicontinuous. Then \eqref{eq:ump} has a solution, the set of solutions
	\[
		x(p,m)\eqdef\argmax_{x\in\budget(p,m)}u(x)
	\]
	is non-empty and compact, and the indirect utility
	$v(p,m)\eqdef\max_{x\in\budget(p,m)}u(x)$ is finite.
\end{proposition}

\begin{proof}
	$\budget(p,m)$ is non-empty, since $0$ belongs to it, and compact by
	\autoref{eg:budget-compact}. Apply \autoref{thm:weierstrass}(ii) and
	\autoref{prop:argmax-compact}(i).
\end{proof}

\paragraph{What the hypotheses mean economically.} Strict positivity of every
price is what bounds the budget set: a free good can be consumed without limit,
and if utility is strictly increasing in it, no optimum exists. Positive wealth
guarantees a non-degenerate problem, though $m=0$ is admissible and yields
$x=0$. Upper semicontinuity of $u$ is weaker than continuity and permits
preferences with an upward jump at a threshold --- a consumer who strictly
prefers reaching a subsistence level to being marginally below it, with the
preferred value attained \emph{at} the threshold. A \emph{downward} jump at the
threshold violates upper semicontinuity, and then the supremum of utility over
the budget set may fail to be attained even though the budget set is compact,
exactly as in \autoref{eg:weierstrass-failures}(i).

\paragraph{What the theorem does not deliver.} Three things.

First, uniqueness. \autoref{prop:demand-exists} produces a non-empty compact
solution \emph{set}. With linear utility $u(x)=\ip{a}{x}$ and prices
proportional to $a$, every point of the budget hyperplane is optimal.
Uniqueness needs strict quasiconcavity (\autoref{ch:convex}).

Second, an interior solution. Nothing in the argument prevents
$x(p,m)$ from lying on the boundary of $\R^n_+$, and indeed corner solutions are
generic when preferences are not strictly monotone in every good. First-order
conditions therefore take the Kuhn--Tucker form of \autoref{ch:optimization}
rather than the gradient-equals-zero form.

Third, continuity of $x(\cdot,\cdot)$ and of $v(\cdot,\cdot)$ in $(p,m)$.
\autoref{prop:demand-exists} is a statement for each fixed $(p,m)$ separately,
and \autoref{eg:separate-vs-joint} is a reminder that pointwise statements do
not aggregate into joint continuity. The budget set itself moves with
$(p,m)$; establishing that the value moves continuously and that the solution
set moves upper hemicontinuously is Berge's maximum theorem
(\autoref{thm:berge-maximum}), which requires the budget correspondence to be
both upper and lower hemicontinuous. That the budget correspondence is
lower hemicontinuous fails at $m=0$, which is the technical reason the
Walrasian demand correspondence is well behaved only for $m>0$.

\begin{remark}[Why upper semicontinuity is the right hypothesis]
	\label{rem:usc-natural}
	Preferences, not utility functions, are the primitive. A preference relation
	$\succsim$ on $\R^n_+$ has \dfnas{closed upper contour sets}{upper contour set}
	if $\{y : y\succsim x\}$ is closed for every $x$; this is a statement about the
	preference relation alone. If $u$ represents $\succsim$, then the upper contour
	set at $x$ is $\{y : u(y)\geq u(x)\}$, so closed upper contour sets are exactly
	upper semicontinuity of $u$ in the sense of \autoref{def:semicontinuity}(ii).
	The hypothesis of \autoref{prop:demand-exists} is therefore a restriction on
	preferences and is invariant to the choice of representing utility function,
	whereas an assumption such as differentiability of $u$ is not.
\end{remark}

\section{Chapter Summary}\label{sec:cont-summary}

Continuity at a point admits three equivalent formulations: the
$\varepsilon$--$\delta$ condition, preservation of limits of sequences
(\autoref{prop:sequential-continuity}), and relative openness of preimages of
open sets (\autoref{thm:continuity-preimage}). The third is metric-free and is
the definition that transfers to general topological spaces and, with
``open'' replaced by ``measurable'', to random variables. Linear maps and norms
are continuous; sums, products, quotients with non-vanishing denominator,
compositions and finite maxima of continuous functions are continuous; separate
continuity in each argument is strictly weaker than joint continuity
(\autoref{eg:separate-vs-joint}). Continuous images of compact sets are compact
and of connected sets are connected, which gives the intermediate value theorem
and the boundedness half of the Weierstrass theorem. Continuity on a compact set
upgrades to uniform continuity (\autoref{thm:heine-cantor}).

Splitting continuity into its two halves isolates what existence of an optimum
actually needs: upper semicontinuity for a maximum, lower semicontinuity for a
minimum, in both cases on a compact domain (\autoref{thm:weierstrass}). The
solution set is then non-empty and compact (\autoref{prop:argmax-compact}). An
arbitrary supremum of lsc functions is lsc, but an arbitrary supremum of
continuous functions need not be usc (\autoref{eg:sup-not-usc}), which is why
the continuity of value functions in their parameters is a separate theorem
rather than a corollary. For the consumer's problem with strictly positive
prices, upper semicontinuity of utility delivers existence of demand and
finiteness of indirect utility, and nothing more.

\begin{exercise}\label{ex:cont-1}
	Prove \autoref{prop:continuity-algebra}(i) directly from
	\autoref{def:continuity}, and identify at which step the hypothesis
	$g(x_0)\neq0$ is used in the quotient case. Then show by example that
	continuity of $f$ and $g$ at $x_0$ is not enough for $f/g$ to have a
	\emph{removable} discontinuity at a zero of $g$.
\end{exercise}

\begin{exercise}\label{ex:cont-2}
	Let $A\subseteq\R^n$ and $f\colon A\to\R$. Show that $f$ is lsc if and only if
	$f=\sup_{i\in I}f_i$ for some family of continuous functions $f_i$ when $A$ is
	compact and $f$ is bounded below. (Hint: for each $x_0$ and each
	$\lambda<f(x_0)$, construct a continuous function below $f$ that takes the value
	$\lambda$ at $x_0$.) Deduce that \autoref{prop:sup-semicontinuous}(i) is sharp.
\end{exercise}

\begin{exercise}\label{ex:cont-3}
	A firm with technology $Y\subseteq\R^n$ closed and prices $p\in\R^n$ solves
	$\pi(p)=\sup_{y\in Y}\ip{p}{y}$. Show that $\pi$ is convex and lsc on $\R^n$ as a
	supremum of linear functions, that $\pi(p)=+\infty$ is possible when $Y$ is
	unbounded, and that if $Y$ is compact then $\pi$ is finite, continuous, and the
	supremum is attained. Relate the failure at unbounded $Y$ to
	\autoref{eg:weierstrass-failures}(ii).
\end{exercise}
